% !TeX root = ../navier-stokes-manuscript.tex
% ============================================================
% 2.5 Simultaneous Sobolev Intersections
% ============================================================

\subsection{Simultaneous Sobolev Intersections}\label{sub:SimultaneousSobolevIntersections}

The ascent has no final spatial height, although each concrete derivative request
uses only finitely many rungs.  To make all \(\partial_x^\alpha u\) with
\(\lvert\alpha\rvert\le k\) continuous, use the three-dimensional embedding
\[
  H^m(\R^3)
  \hookrightarrow
  C^k(\R^3)
\]
whenever
\[
  m>k+\frac32.
\]
Thus each requested order \(k\) chooses one finite height above its threshold.
There is no largest \(m\) as \(k\) ranges over all finite orders, but no single
infinite-order norm is needed.  The proper object is the coordinatewise
intersection
\[
  \bigcap_{m=0}^{\infty}H^m(\R^3),
\]
whose coordinates are all measured on the same field.  For each fixed \(k\),
one coordinate supplies \(C^k\), hence
\[
  \bigcap_{m=0}^{\infty}H^m(\R^3)
  \hookrightarrow
  C^\infty(\R^3).
\]

Time behaves the same way.  Let \(I=(0,T)\), \(T<T_*\), and let \(X\) be a
Hilbert space.
\Needspace{9\baselineskip}
The vector-valued embedding
\[
  H^r(I;X)
  \hookrightarrow
  C^q(\overline I;X)
\]
holds for
\[
  r>q+\frac12.
\]
Each requested time derivative \(q\) chooses one finite \(r\), and all finite
requests give
\[
  \bigcap_{r=0}^{\infty}H^r(I;X)
  \hookrightarrow
  C^\infty(\overline I;X).
\]
No common bound across all \(r\) is asserted.

The mixed jet makes the temporal and spatial choices simultaneously.  For
example, request
\[
  \partial_t^2\partial_x^\alpha u,
  \qquad
  \lvert\alpha\rvert\le4.
\]
The one membership
\[
  u\in H^3\bigl(I;H^6(\R^3)\bigr)
\]
places this entire finite request in
\[
  C^2\bigl(\overline I;C^4(\R^3)\bigr),
\]
since
\[
  3>2+\frac12,
  \qquad
  6>4+\frac32.
\]
The pair \((2,4)\) is only an example.  For arbitrary finite \(q,k\), choose
finite indices with
\[
  r>q+\frac12,
  \qquad
  m>k+\frac32.
\]
Then
\[
  H^r\bigl(I;H^m(\R^3)\bigr)
  \hookrightarrow
  C^q\bigl(\overline I;C^k(\R^3)\bigr).
\]

Collect every finite time--space request on the same velocity by defining
\[
  \mathscr V_\infty(I)
  :=
  \bigcap_{r=0}^{\infty}
  \bigcap_{m=0}^{\infty}
  H^r\bigl(I;H^m(\R^3)\bigr).
\]
Explicitly,
\[
  u\in\mathscr V_\infty(I)
  \quad\Longleftrightarrow\quad
  \norm{u}_{H^r(I;H^m)}
  <
  \infty
  \quad
  \text{for every finite pair }(r,m).
\]
The norms remain separate coordinates; the intersection does not sum them or
replace \(u\) by a family of fields.  The finite embeddings imply
\[
  \mathscr V_\infty(I)
  \hookrightarrow
  C^\infty\bigl(\overline I\times\R^3\bigr).
\]

Choosing one temporal height above the trace threshold and one spatial height
above \(3/2\) places both critical levels on the closed interval:
\[
  u(t)\in H^{1/2}(\R^3)
  \qquad\text{and}\qquad
  u(t)\in H^{3/2}(\R^3)
\]
Their scalar quantities are consequently defined there by
\[
  H(t)=\frac12\norm{\Lambda^{1/2}u(t)}_2^2,
  \qquad
  E_{3/2}(t)=\frac12\norm{\Lambda^{3/2}u(t)}_2^2
\]
Scaling and the first dissipative ascent selected these two heights; the
intersection is what carries them to the terminal face.

Pressure is recovered alongside this velocity space, not inserted into it.  At
finite temporal depth write
\[
  Q_n=\partial_t^np,
\]
The Poisson equation recovers \(Q_n\) from the velocity rows, while the
differentiated momentum equation keeps \(\nabla Q_n\) beside response,
viscosity, and transport.  This produces a pressure-complete jet without
misidentifying pressure as a velocity coordinate.

An actual estimate will control only a finite block.  The two adjacent Sobolev
levels required for a common time trace determine that block’s upper and lower
edges.
